%%==========================================================================
\section{Literature Review}
\label{sec:literature}
%%==========================================================================

The problems studied in the domain of thesis can be divided into the following classes:
\begin{itemize}
	\item Fleet assignment - selection of an aircraft type for each flight;
	\item Aircraft routing - construction of maintenance-feasible sequences for a homogeneous fleet;
	\item Tail assignment - allocation of individual aircraft, with their own locations and maintenance states, to those sequences; and
	\item Maintenance planning - scheduling checks and the tasks performed during them.
\end{itemize}
This distinction matters because models from these streams cannot be compared solely by instance size or solution
time. They differ in decision scope, maintenance semantics, planning horizon,
and treatment of uncertainty \citep{gopalan1998mathematical,
barnhart2003applications, lacasse2010aircraft}.

\subsection{Aircraft routing and tail-assignment formulations}

Early network formulations represent aircraft movement as flow through airport
and time events. The multi-commodity model of \citet{levin1971scheduling}
tracks each aircraft separately, while the time-space network of
\citet{hane1995fleet} provides a natural representation of flight, ground, and
overnight arcs. Maintenance opportunities can then be embedded in feasible
rotations \citep{clarke1996maintenance, clarke1997aircraft}. These models make
physical continuity explicit, but aircraft-specific connection variables grow
rapidly with the number of feasible flight pairs.

Route-based formulations instead associate variables with complete feasible
flight strings. The flight-string model of \citet{barnhart1998flight} exploits
this structure through column generation. Such formulations can express
detailed route feasibility in the pricing problem, but their implementation is
substantially more involved than a directly solved compact MILP. Decomposition
methods offer a related trade-off: they separate coupled decisions to obtain
tractable subproblems, as illustrated by Benders approaches to aircraft routing
and crew scheduling \citep{cordeau2001benders, mercier2005computational}.

Compact formulations avoid explicit route enumeration. The lifted daily
maintenance-routing model of \citet{haouari2012lifted}, for example, strengthens
a compact nonlinear formulation through reformulation and linearisation.
\citet{khaled2018compact} pursue a different compact representation in which
binary variables assign flights directly to individual aircraft and aggregate
time-balance constraints impose route continuity without connection-arc
variables. Its small variable count and direct solvability make it the closest
baseline for this thesis. Compactness, however, transfers more logical work to
aggregate constraints; consequently, corner cases and conditional maintenance
bounds require particularly careful verification.

The daily model of \citet{maher2018daily} is an important operational
counterpoint. It assigns aircraft to one-day lines of flight, enforces current
maintenance needs, and uses look-ahead constraints for the following two days.
Exact and heuristic procedures can modify lines of flight to reduce maintenance
misalignment under day-to-day perturbations. That work addresses operational
flexibility around a supplied route plan, whereas the present thesis audits and
extends a deterministic compact assignment formulation over its full planning
horizon. The two approaches are therefore complementary rather than direct
competitors.

\subsection{Integrated airline planning}

Sequential airline planning can propagate poor or infeasible decisions from
fleet assignment to routing and then to crew planning. \citet{shao2017integrated}
integrate fleet assignment, aircraft routing, and crew pairing in a single
framework and solve it using Benders decomposition with acceleration
strategies. Their results demonstrate the value of preserving cross-stage
dependencies, but also show why decomposition is needed once fleet, route, and
crew decisions are combined.

At a narrower integration level, \citet{mirjafari2020integrated} jointly model
maintenance routing and crew scheduling, including aircraft replacement,
deadheading, night-flight allocation, and flight-hour maintenance thresholds.
They compare Lagrangian relaxation and particle swarm optimisation, with the
relaxation producing solutions closer to the small-instance optimum. These
integrated models establish that routing decisions interact with downstream
resources. They do not, however, provide a like-for-like test of the compact
tail-assignment constraints studied here, because their variables, objectives,
and resource constraints cover a broader planning problem.

\subsection{Robustness and operational uncertainty}

Deterministic cost-optimal routes can be fragile when delays propagate through
tight aircraft connections. Recoverable-robust tail assignment explicitly
balances planned performance with the ability to repair a disrupted schedule
\citep{froyland2013recoverable}. A simulation-optimisation line of work instead
uses stochastic evaluation to place connection buffers and retime flights. The
initial particle-swarm and Monte Carlo framework of
\citet{benahmed2015simulation} was developed into the weekly hybrid model of
\citet{benahmed2017hybrid}. The latter combines a linearised maintenance-routing
model with simulation-based retiming and reports substantial improvements in
on-time performance and propagated delay on airline data.

These studies reveal a limitation of the present thesis: the corrected model
establishes deterministic physical and maintenance feasibility, but it does not
by itself guarantee resilience to stochastic delays. Robust retiming, recourse,
and scenario-based evaluation are therefore complementary extensions, not
features that can be inferred from a feasible deterministic schedule.

\subsection{Maintenance planning at different decision levels}

Aircraft maintenance enters the literature at three levels that should not be
conflated. First, routing models ensure that an aircraft reaches an eligible
station before a usage limit is exceeded. Second, long-term check scheduling
selects dates and capacities for maintenance packages. Third, task allocation
assigns individual inspection and repair tasks to those checks. The present
thesis operates mainly at the first level, while using a stylised hierarchy to
represent interactions among check categories.

At the long-term level, \citet{deng2020maintenance} use dynamic programming to
coordinate A- and C-check schedules for a heterogeneous fleet while accounting
for aircraft status and maintenance capacity. Their airline case study reports
fewer checks and higher aircraft availability than current practice. At the
task level, \citet{witteman2021binpacking} formulate repeated maintenance-task
allocation as a time-constrained variable-sized bin-packing problem. Their
worst-fit-decreasing heuristic obtains small gaps to an exact method while
responding quickly enough for repeated replanning. \citet{deng2021dss} then
integrate check scheduling, task allocation, and shift planning in a decision
support system and demonstrate improvements on airline cases.

This evidence also requires more precise terminology than a simple statement
that four checks are universally mandated. Letter checks are practical packages
whose content and intervals depend on the approved maintenance programme,
aircraft type, utilisation, and operator. The supplied maintenance studies note
that B-check tasks are often absorbed into successive A-checks and that some
operators combine D-check work with heavy C-checks
\citep{deng2020maintenance, witteman2021binpacking, deng2021dss}. Accordingly,
the A/B/C/D structure in Chapter~\ref{ch:hierarchy} is a modelling abstraction
for nested usage limits and counter resets. It is not a claim that every
operator implements four separate maintenance events.

For this thesis, the maintenance-planning literature has two further
implications. Check feasibility depends on both utilisation counters, such as
flight hours and cycles, and calendar limits; and an aircraft's state before the
planning horizon cannot be discarded. Conversely, a routing-level model that
places a check does not solve hangar capacity, task packaging, workforce, or
shift planning. Those resource layers remain outside the present formulation.

\subsection{Heuristic solution methods and cross-domain transfer}

Constructive and local-search methods are operationally useful when exact
solution times are incompatible with replanning. \citet{gronkvist2005tail}
develops tail-assignment heuristics in this setting. The attached locomotive
assignment chapters suggest an additional transferable design: represent
movement and waiting in an acyclic space-time network, select locomotives in a
cost-based order, repeatedly solve a least-cost path problem, update arc costs
after assignment, and terminate a path when a maintenance limit becomes due.
The physical interpretation is railway-specific, but the algorithmic principles
are applicable to aircraft because time also gives the routing network a natural
topological order.

This transfer must preserve aircraft-specific feasibility. Airport continuity,
turnaround time, aircraft eligibility, and all maintenance counters have to be
encoded in nodes, arcs, or state labels before a shortest-path solution is a
valid tail route. Dijkstra's shortest-path method provides the basic path engine
\citep{dijkstra1959note}; ant colony optimisation provides a population-based
alternative for exploring assignment sequences \citep{dorigo1996ant,
dorigo2004ant}. In this thesis these methods are benchmarked as constructive
alternatives to the compact MILP, not treated as evidence of optimality.

\subsection{Research gap and position of the thesis}

The reviewed literature offers sophisticated solutions for route generation,
integrated airline planning, robust retiming, long-term check scheduling, and
maintenance-task allocation. Less attention is paid to auditing the logical
correctness of a compact published formulation before extending or benchmarking
it. This thesis addresses that narrower but foundational gap using
\citet{khaled2018compact} as the direct baseline.

The contribution is positioned in four parts. First, the aggregate ground-flow
conditions are tested against simultaneous and overlapping flight patterns, and
explicit non-overlap conditions are introduced where needed. Second, the
conditional cumulative-hour logic of Constraint~(13) is analysed and replaced
by two endpoint-anchored inequalities. Third, the maintenance state is extended
to multiple nested check categories, including pre-horizon usage, while clearly
remaining at routing level rather than modelling maintenance tasks or workforce.
Fourth, exact and heuristic methods are evaluated on common generated instances
under the same objective and feasibility checks.

This scope supports a controlled claim: the thesis improves the correctness and
maintenance coverage of one compact TAP formulation and evaluates practical
solution alternatives for that formulation. It does not claim to supersede
route-based, decomposition, robust, or maintenance-resource models, whose
broader decision scopes answer different operational questions.
